%%%PAGE 67 rot=0 legibility=good summary=力学三大观点(动力学、能量、动量)综合应用的选用原则：何时用牛二、动量定理/动能定理、两大守恒定律、能量守恒以及碰撞爆炸类问题
%%%BLOCK id=067-01 ch=07 sec=力学三大观点综合应用 kind=method alt=03,04,06 cont=none y=0.07-0.50
\textbf{力学三大观点的综合应用}

直线运动、平抛运动、圆周运动的组合。

动力学、能量、动量的综合应用

要列各物理量在某一时刻的关系式，用牛二

研究某一物体受到力的持续作用发生运动状态改变\quad 用动量定理\ （涉及时间）\ 或动能定理\ （涉及位移） % 原稿“(涉及时间)”“(涉及位移)”写在本行上方，分别位于“用动量定理”“或动能定理”之上

研究系统，用动量守恒定律和机械能守恒定律\quad 注意守恒的条件

涉及相对位移，优先考虑能量守恒定律，系统克服摩擦力所做的总功等于系统机械能的减少量（转变为系统内能的量）。

在涉及碰撞、爆炸、打击、绳绷紧等现象时，一般隐含有系统机械能损失，与其他形式的能量的转换，作用时间都极短，用动量守恒定律解决。
%%%END
%%%PAGEEND 67
